\documentclass[11pt]{article}
\usepackage[a4paper,margin=1in]{geometry}
\usepackage{amsmath,amssymb,amsthm,mathtools}

\newtheorem{theorem}{Theorem}[section]
\newtheorem{lemma}[theorem]{Lemma}
\newtheorem{proposition}[theorem]{Proposition}
\newtheorem{remark}[theorem]{Remark}

\newcommand{\R}{\mathbb R}
\newcommand{\Sph}{\mathbb S^{n-1}}
\newcommand{\Om}{\Omega}
\newcommand{\eps}{\varepsilon}
\newcommand{\cA}{\mathcal A}
\newcommand{\dH}{\,d\mathcal H^{n-1}}
\newcommand{\norm}[1]{\left\|#1\right\|}
\newcommand{\abs}[1]{\left|#1\right|}

\title{A Direct Quantitative Stability Theorem\\
for Small-Amplitude Radial Graphs}
\author{}
\date{May 18, 2026}

\begin{document}
\maketitle

\begin{abstract}
This note proves the sharp Saint--Venant/Fraenkel estimate for a
near-ball radial-graph class without using the BDV theorem as a black box.
The coercive norm is not the full \(H^{1/2}\) norm of the graph.  The
right coercive quantity for the main theorem is the \(L^2\) size of the
radial displacement, which controls the Fraenkel asymmetry.  The proof is
fully quantitative in principle: every constant below is obtained from
elementary Taylor bounds, the Steklov spectrum of the ball, and a
one-dimensional annular Poincare estimate.
\end{abstract}

\section{Statement}

Let
\[
   \Om_R=\{r\theta: \theta\in\Sph,\ 0<r<R(\theta)\},\qquad
   R=1+\varphi,
\]
where \(R\) is smooth and positive.  We normalize
\[
   |\Om_R|=|B_1|,\qquad \int_{\Om_R}x\,dx=0.
\]
Let \(u_\Om\) solve \(-\Delta u_\Om=1\) in \(\Om_R\), \(u_\Om=0\) on
\(\partial\Om_R\), and set \(E(\Om)=-\frac12\int_\Om u_\Om\).

\begin{theorem}[Direct radial-graph stability]\label{thm:direct}
For every \(n\ge2\) there are computable constants
\(\eta_0(n)>0\), \(\eps_0(n)>0\), and \(c(n)>0\) such that the following
holds.  If
\[
   \norm{\varphi}_{L^\infty(\Sph)}\le\eta_0(n),\qquad
   \norm{\varphi}_{L^2(\Sph)}\le\eps_0(n),
\]
then
\[
   E(\Om_R)-E(B_1)\ge c(n)\norm{\varphi}_{L^2(\Sph)}^2.
\]
Consequently
\[
   E(\Om_R)-E(B_1)\ge c_*(n)\,\cA(\Om_R)^2,
\]
with another computable dimensional constant \(c_*(n)>0\).

In particular the theorem applies to any subclass with
\(\norm{\varphi}_{H^{1/2}}\le\eps_0(n)\) and
\(\norm{\varphi}_{L^\infty}\le\eta_0(n)\), since
\(\norm{\varphi}_{L^2}\le\norm{\varphi}_{H^{1/2}}\).
\end{theorem}

\begin{remark}
The \(L^\infty\) amplitude assumption is not a regularity assumption: no
slope, perimeter, curvature, \(C^1\), or \(C^2\) bound appears.  It is used
only to put the boundary in a thin annulus and to make the volume and
barycenter constraints remove the constant and first spherical harmonics
quantitatively.
\end{remark}

\section{The exact energy identity}

Let
\[
   u_0(x)=\frac{1-|x|^2}{2n}.
\]
Define
\[
   z=u_{\Om_R}-u_0 \qquad\hbox{in }\Om_R.
\]
Then \(z\) is harmonic in \(\Om_R\), and on the boundary
\[
   z(R(\theta)\theta)=-u_0(R(\theta)\theta)
      =\frac{R(\theta)^2-1}{2n}
      =:h(\theta).
\]

\begin{lemma}[Exact graph identity]\label{lem:identity}
For every smooth positive radial graph \(R\) with \(|\Om_R|=|B_1|\),
\[
   E(\Om_R)-E(B_1)
   =\frac12\int_{\Om_R}|\nabla z|^2\,dx
     -\frac{1}{2n^2(n+2)}
      \int_{\Sph}\bigl(R^{n+2}-1\bigr)\dH .
\]
\end{lemma}

\begin{proof}
Since \(-\Delta u_0=1\), expanding
\[
   J_{\Om_R}(v)=\frac12\int_{\Om_R}|\nabla v|^2-\int_{\Om_R}v
\]
at \(u_0\) gives
\[
   E(\Om_R)=J_{\Om_R}(u_0)
      +\int_{\partial\Om_R}z\,\partial_\nu u_0\,d\mathcal H^{n-1}
      +\frac12\int_{\Om_R}|\nabla z|^2 .
\]
On \(\partial\Om_R\), \(z=(R^2-1)/(2n)\) and
\(\partial_\nu u_0=-x\cdot\nu/n\).  The radial graph identity
\((x\cdot\nu)\,d\mathcal H^{n-1}=R^n\,d\mathcal H^{n-1}_{\Sph}\) gives
\[
   \int_{\partial\Om_R}z\,\partial_\nu u_0\,d\mathcal H^{n-1}
   =\frac{1}{2n^2}\int_{\Sph}(1-R^2)R^n\dH .
\]
Also
\[
   J_{\Om_R}(u_0)
   =-\frac{|B_1|}{2n}
     +\frac{n+1}{2n^2(n+2)}\int_{\Sph}R^{n+2}\dH ,
\]
using \(|\Om_R|=|B_1|\).  Subtracting the same formula with \(R\equiv1\)
and using \(\int_{\Sph}R^n\,d\mathcal H^{n-1}=|\Sph|\) yields the stated
identity.
\end{proof}

\section{Slope-free annular trace lower bound}

Fix \(0<\eta<1/4\), assume \(1-\eta\le R\le1+\eta\), and put
\(\rho=1-\eta\).  Write the spherical-harmonic expansion
\[
   h=\sum_{k\ge0}h_k,\qquad h_k\in\mathcal H_k.
\]

\begin{lemma}[Annular Steklov lower bound]\label{lem:annular}
There is an explicit number \(b_n(\eta)>0\), with \(b_n(\eta)\to2\) as
\(\eta\downarrow0\), such that every \(w\in H^1(\Om_R)\) whose radial
trace on \(r=R(\theta)\) is \(h(\theta)\) satisfies
\[
   \int_{\Om_R}|\nabla w|^2\,dx
      \ge b_n(\eta)\sum_{k\ge2}\norm{h_k}_{L^2(\Sph)}^2.
\]
One admissible choice is
\[
   b_n(\eta)
   =\frac{2(1-\eta)^{n-2}\mu_n(\eta)}
          {2(1-\eta)^{n-2}+\mu_n(\eta)},
   \qquad
   \mu_n(\eta)=\frac{(1-\eta)^{n-1}}{2\eta}.
\]
\end{lemma}

\begin{proof}
Let \(f(\theta)=w(\rho\theta)\).  The energy inside \(B_\rho\) is bounded
from below by the harmonic extension energy:
\[
   \int_{B_\rho}|\nabla w|^2
      \ge \rho^{n-2}\sum_{k\ge1} k\,\norm{f_k}_{L^2(\Sph)}^2 .
\]
Along each ray, Cauchy's inequality gives
\[
   \int_{\rho}^{R(\theta)} |\partial_r w(r,\theta)|^2 r^{n-1}\,dr
      \ge
   \frac{|h(\theta)-f(\theta)|^2}
        {\int_\rho^{R(\theta)}r^{1-n}\,dr}.
\]
Since \(R(\theta)\le1+\eta\), the denominator is at most
\[
   2\eta(1-\eta)^{1-n}.
\]
Therefore the annular part of the energy is bounded below by
\[
   \mu_n(\eta)\norm{h-f}_{L^2(\Sph)}^2,
   \qquad \mu_n(\eta)=\frac{(1-\eta)^{n-1}}{2\eta}.
\]
For each spherical harmonic degree \(k\ge1\), minimize
\[
   \rho^{n-2}k\,|f_k|^2+\mu_n(\eta)|h_k-f_k|^2
\]
over \(f_k\).  The minimum is
\[
   \frac{\rho^{n-2}k\,\mu_n(\eta)}
        {\rho^{n-2}k+\mu_n(\eta)}
   |h_k|^2 .
\]
This coefficient is increasing in \(k\), so on \(k\ge2\) it is bounded
below by the value at \(k=2\), which is exactly \(b_n(\eta)\).
\end{proof}

\section{Elementary consequences of the constraints}

\begin{lemma}[Removal of neutral modes]\label{lem:modes}
There are computable constants \(C_n\) and \(\eta_1(n)>0\) such that if
\(\norm{\varphi}_{L^\infty}\le\eta_1(n)\), \(|\Om_R|=|B_1|\), and
\(\int_{\Om_R}x\,dx=0\), then
\[
   \norm{P_0\varphi}_{L^2}+\norm{P_1\varphi}_{L^2}
      \le C_n\norm{\varphi}_{L^2}^2 .
\]
Consequently, if also \(\norm{\varphi}_{L^2}\le\eps_1(n)\), then
\[
   \norm{P_{\ge2}\varphi}_{L^2}^2
      \ge \frac{15}{16}\norm{\varphi}_{L^2}^2 .
\]
\end{lemma}

\begin{proof}
The volume condition is
\[
   0=\int_{\Sph}\bigl((1+\varphi)^n-1\bigr)\dH
     =n\int_{\Sph}\varphi\,\dH
       +O_n\!\left(\int_{\Sph}\varphi^2\dH\right),
\]
where the \(O_n\)-constant is explicit for
\(\norm{\varphi}_{L^\infty}\le\eta_1(n)\).  This gives
\(\norm{P_0\varphi}_{L^2}\le C_n\norm{\varphi}_{L^2}^2\).

The barycenter condition is
\[
   0=\int_{\Sph}(1+\varphi)^{n+1}\theta\,\dH
     =(n+1)\int_{\Sph}\varphi\theta\,\dH
       +O_n\!\left(\int_{\Sph}\varphi^2\dH\right),
\]
which gives the \(P_1\) estimate.  The last claim follows by choosing
\(\eps_1(n)\) so that the squared neutral-mode contribution is at most
\(\frac1{16}\norm{\varphi}_{L^2}^2\).
\end{proof}

\begin{lemma}[Moment and boundary data bounds]\label{lem:taylor}
There are computable \(\eta_2(n)>0\) and \(C_n\) such that, if
\(\norm{\varphi}_{L^\infty}\le\eta_2(n)\), then
\[
   \int_{\Sph}\bigl(R^{n+2}-1\bigr)\dH
      \le (n+2)(1+C_n\eta_2)\norm{\varphi}_{L^2}^2
\]
and, for \(h=(R^2-1)/(2n)\),
\[
   \norm{P_{\ge2}h}_{L^2}^2
      \ge \frac{1-C_n\eta_2}{n^2}
            \norm{P_{\ge2}\varphi}_{L^2}^2
         - C_n\norm{\varphi}_{L^2}^4 .
\]
\end{lemma}

\begin{proof}
Taylor's formula gives
\[
   R^{n+2}-1=(n+2)\varphi+\frac{(n+2)(n+1)}2\varphi^2
      +O_n(\eta_2\varphi^2),
\]
and the volume expansion used in Lemma~\ref{lem:modes} gives
\[
   \int_{\Sph}\varphi\,\dH
      =-\frac{n-1}{2}\int_{\Sph}\varphi^2\,\dH
        +O_n(\eta_2\norm{\varphi}_{L^2}^2).
\]
Combining these two identities yields the moment estimate.

For the boundary datum,
\[
   h=\frac{\varphi}{n}+\frac{\varphi^2}{2n}.
\]
Since
\(\norm{P_{\ge2}(\varphi^2)}_{L^2}\le
\norm{\varphi}_{L^\infty}\norm{\varphi}_{L^2}\),
the stated lower bound follows from
\(\norm{a+b}^2\ge (1-\tau)\norm{a}^2-\tau^{-1}\norm{b}^2\), with a fixed
small \(\tau\), and then shrinking \(\eta_2(n)\).
\end{proof}

\section{Proof of the theorem}

Choose \(\eta_0(n)>0\) so small that
\[
   b_n(\eta_0)\ge\frac32,\qquad C_n\eta_0\le\frac1{8},
\]
where \(C_n\) is large enough for Lemmas~\ref{lem:modes} and
\ref{lem:taylor}.  Then choose \(\eps_0(n)>0\) so small that the
quartic terms in Lemmas~\ref{lem:modes} and~\ref{lem:taylor} are
absorbed by \(\frac1{16}\norm{\varphi}_{L^2}^2\).

Apply Lemma~\ref{lem:annular} to \(w=z\), then Lemmas~\ref{lem:modes}
and~\ref{lem:taylor}:
\[
   \int_{\Om_R}|\nabla z|^2
      \ge \frac{3}{2}\norm{P_{\ge2}h}_{L^2}^2
      \ge \frac{5}{4n^2}\norm{\varphi}_{L^2}^2
\]
after the above choices of \(\eta_0,\eps_0\).  Also
\[
   \int_{\Sph}(R^{n+2}-1)\dH
      \le \frac{9}{8}(n+2)\norm{\varphi}_{L^2}^2 .
\]
The exact identity in Lemma~\ref{lem:identity} therefore gives
\[
\begin{aligned}
   E(\Om_R)-E(B_1)
   &\ge
      \frac12\frac{5}{4n^2}\norm{\varphi}_{L^2}^2
      -\frac{1}{2n^2(n+2)}
       \frac{9}{8}(n+2)\norm{\varphi}_{L^2}^2  \\
   &= \frac{1}{16n^2}\norm{\varphi}_{L^2}^2 .
\end{aligned}
\]
This proves the first assertion, with the displayed constant after the
explicit threshold choices above.

Finally,
\[
   |\Om_R\Delta B_1|
      =\frac1n\int_{\Sph}|R^n-1|\dH
      \le C_n\norm{\varphi}_{L^1(\Sph)}
      \le C_n\norm{\varphi}_{L^2(\Sph)}.
\]
Since \(\cA(\Om_R)\le |B_1|^{-1}|\Om_R\Delta B_1|\), the Fraenkel
estimate follows.

\section{What this proves and what it does not prove}

The proof is quantitative and does not invoke BDV's global theorem or
BDV's nearly-spherical theorem.  It proves the main theorem in a genuine
small-amplitude radial-graph class with no slope or curvature control.
The \(H^{1/2}\) hypothesis is used only as a convenient way to guarantee
small \(L^2\); the coercive output is the \(L^2\)/Fraenkel scale.

The proof does not prove the false estimate
\[
   E(\Om_R)-E(B_1)\gtrsim\norm{\varphi}_{H^{1/2}}^2
\]
without additional geometry.  High-frequency capillary ripples still rule
out that stronger statement.  It also does not remove the need for an
amplitude hypothesis: \(H^{1/2}(\Sph)\) alone does not put the graph in a
thin annulus or even give an \(L^\infty\) radius bound.

\end{document}
